%%%%%%%%%%%%%%%%%%%%%%%%%%%%%%%%%%%%%%%%%%%%%%%%%%%%%%%%%%%%%
%% THÉORIE DES CATÉGORIES — Chapitres 3–4
%% Limites et Colimites / Adjonctions
%%%%%%%%%%%%%%%%%%%%%%%%%%%%%%%%%%%%%%%%%%%%%%%%%%%%%%%%%%%%%

%%%%%%%%%%%%%%%%%%%%%%%%%%%%%%%%%%%%%%%%%%%%%%%%%%%%%%%%%%%%
%%           CHAPITRE 3 : LIMITES ET COLIMITES            %%
%%%%%%%%%%%%%%%%%%%%%%%%%%%%%%%%%%%%%%%%%%%%%%%%%%%%%%%%%%%%

\chapter{Limites et Colimites}\label{ch:limites-colimites}
\minitoc

\vspace{1em}

Les limites et colimites sont les généralisations catégoriques de
constructions omniprésentes en mathématiques~: produits, intersections,
noyaux, images réciproques, limites projectives et inductives.
Dans ce chapitre, nous développons la théorie de façon systématique,
en partant du langage des diagrammes et des cônes, puis en énonçant
la propriété universelle qui définit une limite. Nous démontrons
le théorème fondamental d'existence --- les produits et égaliseurs
suffisent pour construire toutes les limites finies --- et nous
établissons la relation cruciale entre limites et foncteurs
représentables.

\index{limite|(}
\index{colimite|(}

%% ============================================================
\section{Diagrammes et catégories d'indices}\label{sec:diagrammes}
%% ============================================================

\begin{definition}[Diagramme]\label{def:diagramme}
\index{diagramme}
\index{catégorie d'indices}
\index{foncteur!diagramme}
Soit $\mathcal{C}$ une catégorie et $\mathcal{J}$ une petite catégorie
(la \emph{catégorie d'indices} ou \emph{catégorie de forme}).
Un \textbf{diagramme de forme~$\mathcal{J}$ dans~$\mathcal{C}$} est un foncteur
\[
    F \colon \mathcal{J} \longrightarrow \mathcal{C}.
\]
On note $F_j$ l'image d'un objet $j \in \Ob(\mathcal{J})$
et $F_\alpha \colon F_j \to F_k$ l'image d'un morphisme
$\alpha \colon j \to k$ dans~$\mathcal{J}$.
\end{definition}

\begin{exemple}\label{ex:formes-diagrammes}
\index{catégorie discrète}
\index{paire parallèle}
\index{span}
\index{cospan}
Les petites catégories suivantes apparaissent constamment comme catégories d'indices.
\begin{enumerate}[label=(\roman*)]
\item \textbf{Catégories discrètes.}
    Soit $\mathcal{J}$ une catégorie discrète sur un ensemble~$I$
    (objets~$I$, seuls morphismes~: les identités).
    Un diagramme $F\colon \mathcal{J}\to\mathcal{C}$ est simplement une famille
    d'objets $\{F_i\}_{i\in I}$.

\item \textbf{La paire parallèle.}
    $\mathcal{J} = (\bullet \rightrightarrows \bullet)$,
    i.e.\ deux objets $0,1$ avec deux morphismes non triviaux
    $\alpha,\beta\colon 0\to 1$.
    Un diagramme de cette forme est une paire de morphismes parallèles
    $f,g\colon A\to B$ dans~$\mathcal{C}$.

\item \textbf{Le span.}
    $\mathcal{J} = (\bullet \leftarrow \bullet \rightarrow \bullet)$.
    Un diagramme est un span $A \xleftarrow{f} C \xrightarrow{g} B$.

\item \textbf{Le cospan.}
    $\mathcal{J} = (\bullet \rightarrow \bullet \leftarrow \bullet)$.
    Un diagramme est un cospan $A \xrightarrow{f} C \xleftarrow{g} B$.
\end{enumerate}
\end{exemple}

\begin{remarque}\label{rem:diagramme-notation}
\index{diagramme!comme foncteur}
Un diagramme n'est rien d'autre qu'un foncteur ; le terme
\og{}diagramme\fg{} signale simplement que l'on s'apprête à en
prendre la limite ou la colimite.
\end{remarque}


%% ============================================================
\section{Cônes et cocônes}\label{sec:cones}
%% ============================================================

\begin{definition}[Cône]\label{def:cone}
\index{cône}
\index{cône!sommet}
\index{cône!jambe}
Soit $F\colon \mathcal{J}\to\mathcal{C}$ un diagramme.
Un \textbf{cône sur~$F$} est une paire $(N,\,\{\pi_j\}_{j\in\mathcal{J}})$
formée d'un objet $N\in\mathcal{C}$ (le \emph{sommet} du cône)
et d'une famille de morphismes $\pi_j\colon N\to F_j$ (les \emph{jambes}),
un pour chaque objet $j\in\mathcal{J}$,
tels que pour tout morphisme $\alpha\colon j\to k$ dans~$\mathcal{J}$
le triangle suivant commute~:
\[
\begin{tikzcd}
& N \arrow[dl, "\pi_j"'] \arrow[dr, "\pi_k"] & \\
F_j \arrow[rr, "F_\alpha"'] & & F_k
\end{tikzcd}
\]
Autrement dit, $F_\alpha \circ \pi_j = \pi_k$ pour tout $\alpha\colon j\to k$.
\end{definition}

\begin{definition}[Morphisme de cônes]\label{def:morphisme-cone}
\index{cône!morphisme}
Soient $(N,\pi_j)$ et $(N',\pi'_j)$ deux cônes sur~$F$.
Un \textbf{morphisme de cônes} de $(N,\pi_j)$ vers $(N',\pi'_j)$
est un morphisme $u\colon N\to N'$ dans~$\mathcal{C}$ tel que
$\pi'_j \circ u = \pi_j$ pour tout $j\in\mathcal{J}$~:
\[
\begin{tikzcd}
N \arrow[rr, "u"] \arrow[dr, "\pi_j"'] && N' \arrow[dl, "\pi'_j"] \\
& F_j &
\end{tikzcd}
\]
Les cônes sur~$F$ et leurs morphismes forment une catégorie,
notée $\mathsf{Cone}(F)$.
\index{cône!catégorie des cônes}
\end{definition}

\begin{definition}[Cocône]\label{def:cocone}
\index{cocône}
\index{cocône!sommet}
\index{cocône!jambe}
Dualement, un \textbf{cocône sous~$F$} est une paire
$(N,\,\{\iota_j\}_{j\in\mathcal{J}})$ où $\iota_j\colon F_j\to N$
et $\iota_k \circ F_\alpha = \iota_j$ pour tout $\alpha\colon j\to k$~:
\[
\begin{tikzcd}
F_j \arrow[rr, "F_\alpha"] \arrow[dr, "\iota_j"'] & & F_k \arrow[dl, "\iota_k"] \\
& N &
\end{tikzcd}
\]
Les cocônes sous~$F$ forment une catégorie $\mathsf{Cocone}(F)$.
\end{definition}

\begin{remarque}\label{rem:cone-transfo-nat}
\index{foncteur constant}
\index{transformation naturelle!et cônes}
Soit $\Delta_N \colon \mathcal{J}\to\mathcal{C}$ le foncteur constant
envoyant tout objet sur~$N$ et tout morphisme sur~$\id_N$.
Alors un cône $(N,\pi_j)$ sur~$F$ est exactement une transformation naturelle
$\pi\colon \Delta_N \Rightarrow F$.
Dualement, un cocône est une transformation naturelle $\iota\colon F\Rightarrow\Delta_N$.
\end{remarque}


%% ============================================================
\section{Limites}\label{sec:limites}
%% ============================================================

\begin{definition}[Limite]\label{def:limite}
\index{limite!définition}
\index{cône universel}
\index{limite!propriété universelle}
Soit $F\colon\mathcal{J}\to\mathcal{C}$ un diagramme.
Une \textbf{limite} de~$F$ est un cône $(\varprojlim F,\,\{\pi_j\}_{j})$
qui est \emph{terminal} dans $\mathsf{Cone}(F)$.
Explicitement~:
\begin{enumerate}[label=(\alph*)]
\item $(\varprojlim F,\,\pi_j)$ est un cône sur~$F$ ;
\item pour tout cône $(N,\,\psi_j)$ sur~$F$, il existe un
      \emph{unique} morphisme $u\colon N\to\varprojlim F$
      tel que $\pi_j\circ u = \psi_j$ pour tout $j\in\mathcal{J}$~:
\end{enumerate}
\[
\begin{tikzcd}
N \arrow[dr, "\psi_j"'] \arrow[rr, dashed, "\exists!\, u"] & & \varprojlim F \arrow[dl, "\pi_j"] \\
& F_j &
\end{tikzcd}
\]
On appelle les $\pi_j$ les \emph{projections canoniques}\index{projection canonique}.
\end{definition}

\begin{proposition}\label{prop:limite-unique}
\index{limite!unicité}
Si elle existe, la limite est unique à unique isomorphisme près.
\end{proposition}

\begin{proof}
Soient $(L,\pi_j)$ et $(L',\pi'_j)$ deux limites de~$F$.
Par la propriété universelle de~$L$, il existe un unique $u\colon L'\to L$
avec $\pi_j\circ u = \pi'_j$. Par celle de~$L'$, il existe un unique
$v\colon L\to L'$ avec $\pi'_j\circ v = \pi_j$.
Alors $u\circ v\colon L\to L$ satisfait $\pi_j\circ(u\circ v)=\pi_j$,
et par unicité $u\circ v = \id_L$. De même $v\circ u = \id_{L'}$.
\end{proof}


%% ============================================================
\section{Colimites}\label{sec:colimites}
%% ============================================================

\begin{definition}[Colimite]\label{def:colimite}
\index{colimite!définition}
\index{cocône universel}
Soit $F\colon\mathcal{J}\to\mathcal{C}$ un diagramme.
Une \textbf{colimite} de~$F$ est un cocône $(\varinjlim F,\,\{\iota_j\}_{j})$
qui est \emph{initial} dans $\mathsf{Cocone}(F)$.
Pour tout cocône $(N,\,\psi_j)$ sous~$F$, il existe un unique
$u\colon \varinjlim F\to N$ tel que $u\circ\iota_j = \psi_j$
pour tout $j$~:
\[
\begin{tikzcd}
& \varinjlim F \arrow[dd, dashed, "\exists!\, u"] & \\
F_j \arrow[ur, "\iota_j"] \arrow[dr, "\psi_j"'] & & F_k \arrow[ul, "\iota_k"'] \arrow[dl, "\psi_k"] \\
& N &
\end{tikzcd}
\]
Les $\iota_j$ sont les \emph{injections canoniques}\index{injection canonique}.
\end{definition}

\begin{remarque}\label{rem:colimite-dualite}
\index{dualité!limite-colimite}
La colimite de $F\colon\mathcal{J}\to\mathcal{C}$ est la limite de
$F^{\op}\colon\mathcal{J}^{\op}\to\mathcal{C}^{\op}$.
Tout énoncé sur les limites se dualise en un énoncé sur les colimites.
\end{remarque}


%% ============================================================
\section{Produits et coproduits}\label{sec:produits-coproduits}
%% ============================================================

\begin{definition}[Produit]\label{def:produit}
\index{produit}
\index{produit!projection}
Soit $\{A_i\}_{i\in I}$ une famille d'objets de~$\mathcal{C}$
(i.e.\ un diagramme de forme discrète~$I$).
Le \textbf{produit} $\prod_{i\in I}A_i$ est la limite de ce diagramme.
Concrètement, c'est un objet $P$ muni de projections
$\pi_i\colon P\to A_i$ tel que, pour tout objet~$X$ et toute famille
de morphismes $f_i\colon X\to A_i$, il existe un unique
$\langle f_i\rangle\colon X\to P$ avec $\pi_i\circ\langle f_i\rangle = f_i$~:
\[
\begin{tikzcd}
& X \arrow[dl, "f_i"'] \arrow[d, dashed, "{\langle f_i\rangle}"] \arrow[dr, "f_j"] & \\
A_i & \prod_{i\in I}A_i \arrow[l, "\pi_i"] \arrow[r, "\pi_j"'] & A_j
\end{tikzcd}
\]
\end{definition}

\begin{definition}[Coproduit]\label{def:coproduit}
\index{coproduit}
\index{coproduit!injection}
Dualement, le \textbf{coproduit} $\coprod_{i\in I}A_i$ est la colimite
du diagramme discret $\{A_i\}_{i\in I}$. C'est un objet~$S$ muni
d'injections $\iota_i\colon A_i\to S$ tel que, pour tout objet~$X$ et
toute famille $f_i\colon A_i\to X$, il existe un unique
$[f_i]\colon S\to X$ avec $[f_i]\circ\iota_i = f_i$.
\end{definition}

\begin{exemple}\label{ex:produits-concrets}
\index{produit!dans $\Set$}
\index{coproduit!dans $\Set$}
\index{produit!dans $\Grp$}
\index{coproduit!dans $\Grp$}
\begin{enumerate}[label=(\roman*)]
\item Dans $\Set$, le produit est le produit cartésien et le coproduit est
      l'union disjointe.
\item Dans $\Grp$, le produit est le produit direct et le coproduit est
      le produit libre.
\item Dans $\Ab$, le produit est le produit direct et le coproduit
      (fini) est la somme directe $\bigoplus_{i\in I}A_i$.
\item Dans $\Top$, le produit est le produit muni de la topologie produit
      et le coproduit est l'union disjointe avec la topologie somme.
\item Dans $\Vect_\K$, produit et coproduit finis coïncident avec
      la somme directe.
\end{enumerate}
\end{exemple}


%% ============================================================
\section{Égaliseurs et coégaliseurs}\label{sec:egaliseurs}
%% ============================================================

\begin{definition}[Égaliseur]\label{def:egaliseur}
\index{égaliseur}
Soient $f,g\colon A\rightrightarrows B$ deux morphismes parallèles.
L'\textbf{égaliseur} de~$f$ et~$g$ est la limite du diagramme
$\mathcal{J}=(\bullet\rightrightarrows\bullet)$, i.e.\ un objet~$E$
muni d'un morphisme $e\colon E\to A$ tel que $f\circ e = g\circ e$,
universel pour cette propriété~:
\[
\begin{tikzcd}
E \arrow[r, "e"] & A \arrow[r, shift left, "f"] \arrow[r, shift right, "g"'] & B
\end{tikzcd}
\]
Pour tout $h\colon X\to A$ avec $f\circ h = g\circ h$,
il existe un unique $\bar{h}\colon X\to E$ avec $e\circ\bar{h}=h$.
\end{definition}

\begin{definition}[Coégaliseur]\label{def:coegaliseur}
\index{coégaliseur}
Le \textbf{coégaliseur} de $f,g\colon A\rightrightarrows B$ est la colimite
du même diagramme~: un objet~$Q$ muni de $q\colon B\to Q$
tel que $q\circ f = q\circ g$, universel pour cette propriété~:
\[
\begin{tikzcd}
A \arrow[r, shift left, "f"] \arrow[r, shift right, "g"'] & B \arrow[r, "q"] & Q
\end{tikzcd}
\]
\end{definition}

\begin{exemple}\label{ex:egaliseurs-concrets}
\index{égaliseur!dans $\Set$}
\index{coégaliseur!dans $\Set$}
\begin{enumerate}[label=(\roman*)]
\item Dans $\Set$, l'égaliseur de $f,g\colon A\rightrightarrows B$ est
      le sous-ensemble $\{a\in A : f(a)=g(a)\}$ avec l'inclusion.
\item Dans $\Set$, le coégaliseur est le quotient $B/{\sim}$ où $\sim$
      est la plus petite relation d'équivalence telle que $f(a)\sim g(a)$
      pour tout $a\in A$.
\item Dans $\Ab$, l'égaliseur de $f,g$ est le noyau de $f-g$.
\end{enumerate}
\end{exemple}


%% ============================================================
\section{Noyaux et conoyaux}\label{sec:noyaux-conoyaux}
%% ============================================================

\begin{definition}[Noyau]\label{def:noyau}
\index{noyau!catégorique}
Dans une catégorie avec objet nul~$0$\index{objet nul}, le \textbf{noyau}
d'un morphisme $f\colon A\to B$ est l'égaliseur de~$f$ et du morphisme
nul $0_{AB}\colon A\to 0\to B$~:
\[
\begin{tikzcd}
\Ker(f) \arrow[r, hookrightarrow, "k"] & A \arrow[r, shift left, "f"] \arrow[r, shift right, "0"'] & B
\end{tikzcd}
\]
\end{definition}

\begin{definition}[Conoyau]\label{def:conoyau}
\index{conoyau}
Le \textbf{conoyau} de $f\colon A\to B$ est le coégaliseur de~$f$ et de~$0_{AB}$~:
\[
\begin{tikzcd}
A \arrow[r, shift left, "f"] \arrow[r, shift right, "0"'] & B \arrow[r, twoheadrightarrow, "c"] & \coker(f)
\end{tikzcd}
\]
\end{definition}

\begin{exemple}\label{ex:noyaux-concrets}
\index{noyau!dans $\Ab$}
\index{noyau!dans $\Mod$}
Dans $\Ab$ ou $\Mod_R$, le noyau et le conoyau catégoriques coïncident
avec le noyau et le conoyau algébriques habituels.
\end{exemple}


%% ============================================================
\section{Produits fibrés et sommes amalgamées}\label{sec:pullbacks-pushouts}
%% ============================================================

\begin{definition}[Produit fibré]\label{def:produit-fibre}
\index{produit fibré}
\index{pullback|see{produit fibré}}
Soit un cospan $A \xrightarrow{f} C \xleftarrow{g} B$.
Le \textbf{produit fibré} (ou \emph{pullback}) est la limite de ce diagramme.
C'est un objet $A\times_C B$ muni de $p\colon A\times_C B\to A$ et
$q\colon A\times_C B\to B$ tels que $f\circ p = g\circ q$,
universel pour cette propriété~:
\[
\begin{tikzcd}
A\times_C B \arrow[r, "q"] \arrow[d, "p"'] \arrow[dr, phantom, "\lrcorner", very near start]
    & B \arrow[d, "g"] \\
A \arrow[r, "f"'] & C
\end{tikzcd}
\]
Le carré ci-dessus est appelé \emph{carré cartésien}\index{carré cartésien}.
\end{definition}

\begin{definition}[Somme amalgamée]\label{def:somme-amalgamee}
\index{somme amalgamée}
\index{pushout|see{somme amalgamée}}
Soit un span $A \xleftarrow{f} C \xrightarrow{g} B$.
La \textbf{somme amalgamée} (ou \emph{pushout}) est la colimite de ce diagramme~:
un objet $A\sqcup_C B$ muni de $i\colon A\to A\sqcup_C B$ et
$j\colon B\to A\sqcup_C B$ tels que $i\circ f = j\circ g$~:
\[
\begin{tikzcd}
C \arrow[r, "g"] \arrow[d, "f"'] & B \arrow[d, "j"] \\
A \arrow[r, "i"'] & A\sqcup_C B \arrow[ul, phantom, "\ulcorner", very near start]
\end{tikzcd}
\]
\end{definition}

\begin{exemple}\label{ex:pullback-concrets}
\index{produit fibré!dans $\Set$}
\index{somme amalgamée!dans $\Set$}
\index{produit fibré!dans $\Top$}
\begin{enumerate}[label=(\roman*)]
\item Dans $\Set$, $A\times_C B = \{(a,b)\in A\times B : f(a)=g(b)\}$.
\item Dans $\Set$, $A\sqcup_C B = (A\sqcup B)/{\sim}$ où $f(c)\sim g(c)$
      pour tout $c\in C$.
\item Dans $\Top$, le produit fibré est le sous-espace correspondant du produit.
\item Dans $\Grp$, la somme amalgamée est le produit amalgamé de groupes.
\end{enumerate}
\end{exemple}

\begin{proposition}\label{prop:pullback-egaliseur-produit}
\index{produit fibré!via produit et égaliseur}
Si $\mathcal{C}$ admet des produits binaires et des égaliseurs, alors
$\mathcal{C}$ admet des produits fibrés. Plus précisément, le produit fibré
de $A\xrightarrow{f}C\xleftarrow{g}B$ est l'égaliseur de
\[
\begin{tikzcd}
A\times B \arrow[r, shift left, "f\circ\pi_A"] \arrow[r, shift right, "g\circ\pi_B"'] & C
\end{tikzcd}
\]
\end{proposition}

\begin{proof}
Soit $e\colon E\to A\times B$ l'égaliseur de $f\circ\pi_A$ et $g\circ\pi_B$.
Posons $p=\pi_A\circ e$ et $q=\pi_B\circ e$.
Alors $f\circ p = f\circ\pi_A\circ e = g\circ\pi_B\circ e = g\circ q$,
donc $(E,p,q)$ est un cône sur le cospan.
Soit $(X,p',q')$ un autre cône, i.e.\ $f\circ p'=g\circ q'$.
Par la propriété universelle du produit, il existe un unique
$h=\langle p',q'\rangle\colon X\to A\times B$ avec
$\pi_A\circ h=p'$ et $\pi_B\circ h=q'$.
Alors $f\circ\pi_A\circ h = f\circ p' = g\circ q' = g\circ\pi_B\circ h$,
donc par la propriété universelle de l'égaliseur, $h$ se factorise
uniquement par~$e$.
\end{proof}


%% ============================================================
\section{Théorème d'existence des limites finies}\label{sec:existence-limites}
%% ============================================================

\begin{theoreme}\label{thm:existence-limites-finies}
\index{limite!théorème d'existence}
\index{catégorie finiment complète}
Soit $\mathcal{C}$ une catégorie.
Les conditions suivantes sont équivalentes~:
\begin{enumerate}[label=(\roman*)]
\item $\mathcal{C}$ admet toutes les limites finies (i.e.\ les limites
      de diagrammes $F\colon\mathcal{J}\to\mathcal{C}$ avec $\mathcal{J}$ finie).
\item $\mathcal{C}$ admet les produits finis et les égaliseurs.
\item $\mathcal{C}$ admet un objet terminal, les produits binaires et
      les égaliseurs.
\end{enumerate}
On dit alors que $\mathcal{C}$ est \textbf{finiment complète}\index{finiment complète}.
\end{theoreme}

\begin{proof}
L'implication (i)$\Rightarrow$(ii) est claire puisque les produits finis
et les égaliseurs sont des limites finies.

(ii)$\Rightarrow$(iii)~: le produit vide est un objet terminal.

(iii)$\Rightarrow$(i)~: soit $F\colon\mathcal{J}\to\mathcal{C}$ un diagramme
avec $\mathcal{J}$ finie, disons $\Ob(\mathcal{J})=\{j_1,\dots,j_n\}$
et $\Mor(\mathcal{J})$ (non identités) $= \{\alpha_1,\dots,\alpha_m\}$
avec $\alpha_k\colon s(k)\to t(k)$.
Formons le produit $P = \prod_{i=1}^n F_{j_i}$ avec projections $\pi_i$,
et le produit $Q = \prod_{k=1}^m F_{t(k)}$ avec projections $\rho_k$.

Définissons deux morphismes $\varphi,\psi\colon P\to Q$ par~:
\[
    \rho_k\circ\varphi = F_{\alpha_k}\circ\pi_{s(k)},
    \qquad
    \rho_k\circ\psi = \pi_{t(k)}.
\]
Ces morphismes existent par la propriété universelle de~$Q$.
Soit $e\colon E\to P$ l'égaliseur de $\varphi$ et~$\psi$.

Montrons que $(E, \{\pi_i\circ e\}_i)$ est une limite de~$F$.
Pour tout $\alpha_k\colon s(k)\to t(k)$~:
\[
    F_{\alpha_k}\circ(\pi_{s(k)}\circ e)
    = \rho_k\circ\varphi\circ e
    = \rho_k\circ\psi\circ e
    = \pi_{t(k)}\circ e,
\]
donc $(E, \pi_i\circ e)$ est bien un cône.
Soit $(N,\psi_i)$ un autre cône. La famille $\psi_i$ induit un unique
$h\colon N\to P$ avec $\pi_i\circ h = \psi_i$.
La condition de cône donne $\varphi\circ h = \psi\circ h$, donc
$h$ se factorise uniquement par~$e$, ce qui conclut.
\end{proof}

\begin{corollaire}\label{cor:existence-colimites-finies}
\index{catégorie finiment cocomplète}
Dualement, $\mathcal{C}$ admet toutes les colimites finies
si et seulement si elle admet les coproduits finis et les coégaliseurs.
On dit alors que $\mathcal{C}$ est \textbf{finiment cocomplète}.
\end{corollaire}


%% ============================================================
\section{$\Hom$ préserve les limites}\label{sec:hom-preserve-limites}
%% ============================================================

\begin{definition}[Foncteur qui préserve les limites]\label{def:preservation-limites}
\index{foncteur!qui préserve les limites}
\index{foncteur continu}
Un foncteur $G\colon\mathcal{C}\to\mathcal{D}$ \textbf{préserve les limites}
de forme~$\mathcal{J}$ si, pour tout diagramme $F\colon\mathcal{J}\to\mathcal{C}$
dont la limite $(L,\pi_j)$ existe, le cône $(GL, G\pi_j)$ est une limite de
$GF\colon\mathcal{J}\to\mathcal{D}$.
Un foncteur qui préserve toutes les petites limites est dit \textbf{continu}\index{continu}.
\end{definition}

\begin{theoreme}\label{thm:hom-preserve-limites}
\index{Hom!préserve les limites}
\index{foncteur représentable!préserve les limites}
Pour tout objet $X\in\mathcal{C}$, le foncteur covariant
\[
    \Hom_{\mathcal{C}}(X,-)\colon\mathcal{C}\longrightarrow\Set
\]
préserve toutes les limites qui existent dans~$\mathcal{C}$.
Autrement dit, si $(L,\pi_j)=\varprojlim F$, alors
\[
    \Hom(X,\varprojlim F) \cong \varprojlim_j\, \Hom(X,F_j)
\]
naturellement en~$X$.
\end{theoreme}

\begin{proof}
Soit $F\colon\mathcal{J}\to\mathcal{C}$ un diagramme avec limite $(L,\pi_j)$.
La limite $\varprojlim_j\Hom(X,F_j)$ dans $\Set$ est
\[
    \bigl\{(h_j)_{j\in\mathcal{J}}\in\textstyle\prod_j\Hom(X,F_j)
    : F_\alpha\circ h_j = h_k \text{ pour tout } \alpha\colon j\to k\bigr\},
\]
c'est-à-dire l'ensemble des cônes de sommet~$X$ sur~$F$.

Définissons $\Phi\colon\Hom(X,L)\to\varprojlim_j\Hom(X,F_j)$ par
$\Phi(u)=(\pi_j\circ u)_j$. Pour tout $u$, la famille $(\pi_j\circ u)_j$
est un cône (car $(\pi_j)$ l'est), donc $\Phi$ est bien définie.

\emph{Injectivité}~: si $\Phi(u)=\Phi(v)$, alors $\pi_j\circ u=\pi_j\circ v$
pour tout~$j$, et par unicité dans la propriété universelle de~$L$, $u=v$.

\emph{Surjectivité}~: soit $(h_j)_j$ un cône de sommet~$X$.
Par la propriété universelle de~$L$, il existe un unique $u\colon X\to L$
avec $\pi_j\circ u = h_j$, i.e.\ $\Phi(u)=(h_j)_j$.

La naturalité en~$X$ est immédiate.
\end{proof}

\begin{corollaire}\label{cor:hom-contravariant-colimites}
\index{Hom!contravariant et colimites}
Le foncteur contravariant $\Hom_{\mathcal{C}}(-,Y)\colon\mathcal{C}^{\op}\to\Set$
transforme les colimites de~$\mathcal{C}$ en limites de~$\Set$~:
\[
    \Hom(\varinjlim F_j,\, Y) \cong \varprojlim_j\, \Hom(F_j,Y).
\]
\end{corollaire}


%% ============================================================
\section{Limites filtrantes}\label{sec:limites-filtrantes}
%% ============================================================

\begin{definition}[Catégorie filtrante]\label{def:categorie-filtrante}
\index{catégorie filtrante}
\index{filtrant}
Une petite catégorie $\mathcal{J}$ est \textbf{filtrante} si~:
\begin{enumerate}[label=(\alph*)]
\item $\mathcal{J}$ est non vide ;
\item pour tous $j,k\in\mathcal{J}$, il existe un objet~$\ell$ et des
      morphismes $j\to\ell$ et $k\to\ell$ ;
\item pour tous morphismes parallèles $f,g\colon j\rightrightarrows k$,
      il existe $h\colon k\to\ell$ tel que $h\circ f = h\circ g$.
\end{enumerate}
\end{definition}

\begin{definition}[Colimite filtrante]\label{def:colimite-filtrante}
\index{colimite!filtrante}
Une \textbf{colimite filtrante} est une colimite de diagramme
$F\colon\mathcal{J}\to\mathcal{C}$ où $\mathcal{J}$ est filtrante.
Dualement, une \textbf{limite cofiltrante}\index{limite!cofiltrante}
est une limite indexée par une catégorie cofiltrante
($\mathcal{J}^{\op}$ filtrante).
\end{definition}

\begin{exemple}\label{ex:colimites-filtrantes}
\index{système inductif}
\begin{enumerate}[label=(\roman*)]
\item Tout ensemble ordonné filtrant $(I,\leq)$ (vu comme catégorie)
      est filtrant. Les colimites indexées par un tel ensemble sont
      les \emph{limites inductives} classiques.
\item $\varinjlim_{n\in\N}\Z/p^n\Z \cong \Z[1/p]/\Z$ dans $\Ab$.
\end{enumerate}
\end{exemple}

\begin{theoreme}\label{thm:colimites-filtrantes-set}
\index{colimite!filtrante dans $\Set$}
Dans $\Set$, les colimites filtrantes commutent avec les limites finies.
Plus précisément, si $\mathcal{J}$ est filtrante et $\mathcal{K}$ est finie, alors
pour $F\colon\mathcal{J}\times\mathcal{K}\to\Set$~:
\[
    \varinjlim_j\, \varprojlim_k\, F(j,k)
    \cong \varprojlim_k\, \varinjlim_j\, F(j,k).
\]
\end{theoreme}

\begin{proof}[Esquisse de démonstration]
On vérifie directement que la description ensembliste de la colimite
filtrante (classes d'équivalence de paires $(j,x_j)$) permet de
commuter les deux constructions, en utilisant les propriétés (a)--(c)
de la catégorie filtrante pour traiter les cas des produits finis
et des égaliseurs.
\end{proof}


%% ============================================================
\section{Exercices}\label{sec:exo-limites}
%% ============================================================

\begin{exercice}\label{exo:produit-fibre-mono}
\index{monomorphisme!et produit fibré}
Soient $f\colon A\to C$ et $g\colon B\to C$ dans une catégorie
avec produits fibrés. Montrer que si $g$ est un monomorphisme,
alors la projection $p\colon A\times_C B\to A$ est un monomorphisme.
\end{exercice}

\begin{exercice}\label{exo:egaliseur-mono}
\index{égaliseur!est un monomorphisme}
Montrer que tout égaliseur est un monomorphisme.
\end{exercice}

\begin{exercice}\label{exo:limite-fonctorielle}
Soient $F,G\colon\mathcal{J}\to\mathcal{C}$ deux diagrammes et
$\alpha\colon F\Rightarrow G$ une transformation naturelle.
Si $\varprojlim F$ et $\varprojlim G$ existent, montrer qu'il
existe un unique morphisme
$\varprojlim\alpha\colon\varprojlim F\to\varprojlim G$
compatible avec les projections.
\end{exercice}

\begin{exercice}\label{exo:produit-dans-cat}
\index{produit!dans $\Cat$}
Décrire explicitement le produit de deux petites catégories
$\mathcal{C}\times\mathcal{D}$ dans $\Cat$.
\end{exercice}

\begin{exercice}\label{exo:limites-dans-foncteurs}
\index{limite!dans une catégorie de foncteurs}
Soit $\mathcal{C}$ une catégorie et $\mathcal{D}$ une catégorie
admettant toutes les limites de forme~$\mathcal{J}$.
Montrer que $\Fun(\mathcal{C},\mathcal{D})$ admet aussi les limites
de forme~$\mathcal{J}$, calculées \og{}objectivement\fg{}~:
$(\varprojlim F_i)(C) = \varprojlim F_i(C)$.
\end{exercice}

\begin{exercice}\label{exo:complete-cocomplete}
\index{catégorie complète}
Montrer que $\Set$, $\Grp$, $\Ab$, $\Top$ et $\Mod_R$
sont des catégories complètes et cocomplètes
(admettent toutes les petites limites et colimites).
\end{exercice}

\begin{exercice}\label{exo:colimite-filtrante-exacte}
\index{colimite!filtrante!exactitude}
Montrer que dans $\Ab$, les colimites filtrantes sont exactes,
i.e.\ le foncteur $\varinjlim\colon\Fun(\mathcal{J},\Ab)\to\Ab$
(pour $\mathcal{J}$ filtrante) est un foncteur exact.
\end{exercice}

\begin{exercice}\label{exo:pushout-dans-grp}
\index{somme amalgamée!dans $\Grp$}
Soient $H\xrightarrow{f}G_1$ et $H\xrightarrow{g}G_2$ des
homomorphismes de groupes. Montrer que la somme amalgamée
$G_1\sqcup_H G_2$ dans $\Grp$ est le produit amalgamé
$G_1 *_H G_2$.
\end{exercice}

\index{limite|)}
\index{colimite|)}


%%%%%%%%%%%%%%%%%%%%%%%%%%%%%%%%%%%%%%%%%%%%%%%%%%%%%%%%%%%%
%%           CHAPITRE 4 : ADJONCTIONS                     %%
%%%%%%%%%%%%%%%%%%%%%%%%%%%%%%%%%%%%%%%%%%%%%%%%%%%%%%%%%%%%

\chapter{Adjonctions}\label{ch:adjonctions}
\minitoc

\vspace{1em}

Les adjonctions sont l'un des concepts les plus fondamentaux de la théorie
des catégories, reliant deux foncteurs de directions opposées par une
bijection naturelle entre ensembles de morphismes.
Introduites par Daniel Kan en 1958\index{Kan, Daniel},
elles unifient une multitude de constructions mathématiques~:
les foncteurs libres, la tensorisation, l'image directe et réciproque
de faisceaux, les suspensions et espaces de lacets.
Comme le disait Saunders Mac~Lane, \og{}les adjonctions
apparaissent partout\fg{}.

\index{adjonction|(}

%% ============================================================
\section{Définition par bijection naturelle}\label{sec:adjonction-def}
%% ============================================================

\begin{definition}[Adjonction]\label{def:adjonction}
\index{adjonction!définition}
\index{adjoint à gauche}
\index{adjoint à droite}
\index{$F\dashv G$}
Soient $\mathcal{C}$ et $\mathcal{D}$ deux catégories et
\[
\begin{tikzcd}
\mathcal{C} \arrow[r, bend left=20, "F", ""{name=A,below}]
& \mathcal{D} \arrow[l, bend left=20, "G", ""{name=B,above}]
\arrow[from=A, to=B, symbol=\dashv]
\end{tikzcd}
\]
deux foncteurs. On dit que $F$ est \textbf{adjoint à gauche} de~$G$
(et $G$ \textbf{adjoint à droite} de~$F$), noté $F\dashv G$,
s'il existe une bijection
\[
    \Phi_{A,B}\colon \Hom_{\mathcal{D}}(FA,\, B)
    \xrightarrow{\;\sim\;}
    \Hom_{\mathcal{C}}(A,\, GB)
\]
naturelle en $A\in\mathcal{C}$ et $B\in\mathcal{D}$.
La naturalité signifie que pour tous $f\colon A'\to A$ dans $\mathcal{C}$
et $g\colon B\to B'$ dans $\mathcal{D}$, le diagramme suivant commute~:
\[
\begin{tikzcd}
\Hom(FA,B) \arrow[r, "\Phi_{A,B}"] \arrow[d, "{(Ff)^*\circ g_*}"']
& \Hom(A,GB) \arrow[d, "{f^*\circ (Gg)_*}"] \\
\Hom(FA',B') \arrow[r, "\Phi_{A',B'}"']
& \Hom(A',GB')
\end{tikzcd}
\]
\end{definition}

\begin{notation}\label{not:adjonction-transposee}
\index{transposée}
\index{morphisme adjoint}
Si $h\colon FA\to B$, on appelle $\tilde{h}=\Phi(h)\colon A\to GB$
la \emph{transposée} (ou \emph{adjointe}) de~$h$. Réciproquement,
si $k\colon A\to GB$, on note $\hat{k}=\Phi^{-1}(k)\colon FA\to B$.
\end{notation}


%% ============================================================
\section{Unité et counité}\label{sec:unite-counite}
%% ============================================================

\begin{definition}[Unité et counité]\label{def:unite-counite}
\index{unité!d'une adjonction}
\index{counité!d'une adjonction}
\index{$\eta$ (unité)}
\index{$\varepsilon$ (counité)}
Soit $F\dashv G$ une adjonction avec bijection~$\Phi$.
\begin{enumerate}[label=(\roman*)]
\item L'\textbf{unité} est la transformation naturelle
$\eta\colon\Id_{\mathcal{C}}\Rightarrow GF$ définie par
\[
    \eta_A = \Phi_{A,FA}(\id_{FA})\colon A\to GFA.
\]
\item La \textbf{counité} est la transformation naturelle
$\varepsilon\colon FG\Rightarrow\Id_{\mathcal{D}}$ définie par
\[
    \varepsilon_B = \Phi^{-1}_{GB,B}(\id_{GB})\colon FGB\to B.
\]
\end{enumerate}
\end{definition}

\begin{proposition}\label{prop:transposee-via-unite}
\index{transposée!via unité et counité}
Pour tout $h\colon FA\to B$ et tout $k\colon A\to GB$~:
\[
    \tilde{h} = Gh\circ\eta_A, \qquad \hat{k} = \varepsilon_B\circ Fk.
\]
\end{proposition}

\begin{proof}
Par naturalité de~$\Phi$ en~$B$, pour $h\colon FA\to B$~:
\[
    \Phi_{A,B}(h)
    = \Phi_{A,B}(h\circ\id_{FA})
    = Gh \circ \Phi_{A,FA}(\id_{FA})
    = Gh\circ\eta_A.
\]
De même, par naturalité en~$A$, pour $k\colon A\to GB$~:
\[
    \Phi^{-1}_{A,B}(k)
    = \Phi^{-1}_{A,B}(\id_{GB}\circ k)
    = \Phi^{-1}_{GB,B}(\id_{GB})\circ Fk
    = \varepsilon_B\circ Fk. \qedhere
\]
\end{proof}


%% ============================================================
\section{Identités triangulaires}\label{sec:identites-triangulaires}
%% ============================================================

\begin{theoreme}[Identités triangulaires]\label{thm:identites-triangulaires}
\index{identités triangulaires}
\index{identités zigzag|see{identités triangulaires}}
Soit $F\dashv G$ avec unité~$\eta$ et counité~$\varepsilon$.
Les compositions suivantes sont des identités~:
\begin{enumerate}[label=(\roman*)]
\item $\varepsilon F\circ F\eta = \id_F$, i.e.\ pour tout $A$~:
\[
\begin{tikzcd}
FA \arrow[r, "F\eta_A"] \arrow[dr, "\id_{FA}"'] & FGFA \arrow[d, "\varepsilon_{FA}"] \\
& FA
\end{tikzcd}
\]
\item $G\varepsilon\circ\eta G = \id_G$, i.e.\ pour tout $B$~:
\[
\begin{tikzcd}
GB \arrow[r, "\eta_{GB}"] \arrow[dr, "\id_{GB}"'] & GFGB \arrow[d, "G\varepsilon_B"] \\
& GB
\end{tikzcd}
\]
\end{enumerate}
\end{theoreme}

\begin{proof}
(i) Soit $A\in\mathcal{C}$. En utilisant la \autoref{prop:transposee-via-unite}
avec $h=\id_{FA}\colon FA\to FA$~:
\[
    \tilde{\id}_{FA} = G(\id_{FA})\circ\eta_A = \eta_A.
\]
Maintenant, appliquons $\Phi^{-1}$ à $\eta_A = \tilde{\id}_{FA}$~:
\[
    \hat{\eta}_A = \varepsilon_{FA}\circ F\eta_A.
\]
Mais $\hat{\eta}_A = \Phi^{-1}(\Phi(\id_{FA})) = \id_{FA}$,
donc $\varepsilon_{FA}\circ F\eta_A = \id_{FA}$.

(ii) Soit $B\in\mathcal{D}$. Avec $k=\id_{GB}\colon GB\to GB$~:
\[
    \hat{\id}_{GB} = \varepsilon_B\circ F(\id_{GB}) = \varepsilon_B.
\]
Appliquons $\Phi$ à $\varepsilon_B = \hat{\id}_{GB}$~:
\[
    \tilde{\varepsilon}_B = G\varepsilon_B\circ\eta_{GB}.
\]
Mais $\tilde{\varepsilon}_B = \Phi(\Phi^{-1}(\id_{GB})) = \id_{GB}$,
donc $G\varepsilon_B\circ\eta_{GB} = \id_{GB}$.
\end{proof}


%% ============================================================
\section{Équivalence des définitions}\label{sec:equivalence-def-adjonctions}
%% ============================================================

\begin{theoreme}\label{thm:equivalence-adjonctions}
\index{adjonction!équivalence des définitions}
Les données suivantes sont équivalentes~:
\begin{enumerate}[label=(\alph*)]
\item Une adjonction $F\dashv G$ donnée par une bijection naturelle
      $\Phi\colon\Hom(FA,B)\xrightarrow{\sim}\Hom(A,GB)$.
\item Un quadruplet $(F,G,\eta,\varepsilon)$ où
      $\eta\colon\Id_\mathcal{C}\Rightarrow GF$ et
      $\varepsilon\colon FG\Rightarrow\Id_\mathcal{D}$
      sont des transformations naturelles satisfaisant les
      identités triangulaires.
\end{enumerate}
\end{theoreme}

\begin{proof}
\emph{(a)$\Rightarrow$(b)}~: c'est ce que nous avons construit dans les
\autoref{def:unite-counite} et \autoref{thm:identites-triangulaires}.

\emph{(b)$\Rightarrow$(a)}~: étant donné $(F,G,\eta,\varepsilon)$
satisfaisant les identités triangulaires, définissons
\[
    \Phi_{A,B}(h) = Gh\circ\eta_A
    \quad\text{pour}\quad h\colon FA\to B,
\]
\[
    \Psi_{A,B}(k) = \varepsilon_B\circ Fk
    \quad\text{pour}\quad k\colon A\to GB.
\]

Vérifions que $\Psi\circ\Phi = \id$.
Soit $h\colon FA\to B$~:
\[
    \Psi(\Phi(h))
    = \varepsilon_B\circ F(Gh\circ\eta_A)
    = \varepsilon_B\circ FGh\circ F\eta_A.
\]
Par naturalité de~$\varepsilon$ (appliquée à $h\colon FA\to B$)~:
$\varepsilon_B\circ FGh = h\circ\varepsilon_{FA}$, donc
\[
    \Psi(\Phi(h))
    = h\circ\varepsilon_{FA}\circ F\eta_A
    = h\circ\id_{FA} = h,
\]
en utilisant la première identité triangulaire.

Vérifions que $\Phi\circ\Psi=\id$.
Soit $k\colon A\to GB$~:
\[
    \Phi(\Psi(k))
    = G(\varepsilon_B\circ Fk)\circ\eta_A
    = G\varepsilon_B\circ GFk\circ\eta_A.
\]
Par naturalité de~$\eta$ (appliquée à $k\colon A\to GB$)~:
$GFk\circ\eta_A = \eta_{GB}\circ k$, donc
\[
    \Phi(\Psi(k))
    = G\varepsilon_B\circ\eta_{GB}\circ k
    = \id_{GB}\circ k = k,
\]
en utilisant la seconde identité triangulaire.

La naturalité de~$\Phi$ en~$A$ et~$B$ suit de la fonctorialité de~$G$
et de la naturalité de~$\eta$.
\end{proof}


%% ============================================================
\section{Exemples fondamentaux}\label{sec:exemples-adjonctions}
%% ============================================================

\begin{exemple}[Libre $\dashv$ Oubli]\label{ex:libre-oubli}
\index{adjonction!libre-oubli}
\index{foncteur libre}
\index{foncteur d'oubli}
Le foncteur d'oubli $U\colon\Grp\to\Set$ admet un adjoint à gauche $F$,
le foncteur \og{}groupe libre\fg{}~:
\[
\begin{tikzcd}
\Set \arrow[r, bend left=20, "F", ""{name=A,below}]
& \Grp \arrow[l, bend left=20, "U", ""{name=B,above}]
\arrow[from=A, to=B, symbol=\dashv]
\end{tikzcd}
\qquad
\Hom_{\Grp}(FX,\,G) \cong \Hom_{\Set}(X,\,UG).
\]
L'unité $\eta_X\colon X\to UFX$ est l'inclusion de~$X$ comme
générateurs dans le groupe libre $FX$. La counité
$\varepsilon_G\colon FUG\to G$ envoie chaque mot formel dans
le groupe libre sur les éléments de~$G$ sur son évaluation dans~$G$.

De même pour $\Ab$, $\Ring$, $\Mod_R$, $\Vect_\K$, etc.
\end{exemple}

\begin{exemple}[Produit $\dashv$ Hom interne]\label{ex:produit-hom}
\index{adjonction!produit-hom}
\index{catégorie cartésienne fermée}
\index{exponentiation}
Dans $\Set$ (ou plus généralement dans toute catégorie cartésienne fermée),
pour tout ensemble~$A$~:
\[
    (-)\times A \;\dashv\; (-)^A,
\]
i.e.\ $\Hom(X\times A,\, B)\cong\Hom(X,\, B^A)$ où $B^A=\Hom(A,B)$.
C'est la \emph{curryfication}\index{curryfication}.
\end{exemple}

\begin{exemple}[Image réciproque $\dashv$ image directe]\label{ex:faisceaux}
\index{adjonction!$f^*\dashv f_*$}
\index{image directe}
\index{image réciproque}
Soit $f\colon X\to Y$ une application continue entre espaces topologiques.
Les foncteurs image réciproque $f^*$ et image directe $f_*$ sur les
faisceaux satisfont $f^*\dashv f_*$~:
\[
    \Hom_{\mathsf{Sh}(X)}(f^*\mathcal{G},\,\mathcal{F})
    \cong \Hom_{\mathsf{Sh}(Y)}(\mathcal{G},\,f_*\mathcal{F}).
\]
\end{exemple}

\begin{exemple}[Suspension $\dashv$ Lacets]\label{ex:suspension-lacets}
\index{adjonction!$\Sigma\dashv\Omega$}
\index{suspension}
\index{espace de lacets}
Dans la catégorie homotopique pointée $\Top_*$, la suspension réduite
$\Sigma$ et le foncteur espace de lacets~$\Omega$ satisfont
\[
    \Sigma \dashv \Omega,\qquad
    [\Sigma X,\,Y]_* \cong [X,\,\Omega Y]_*
\]
où $[-,-]_*$ désigne les classes d'homotopie pointée.
\end{exemple}

\begin{exemple}[Foncteur diagonal $\dashv$ Produit]\label{ex:diagonal-produit}
\index{adjonction!diagonal-produit}
\index{foncteur diagonal}
Soit $\Delta\colon\mathcal{C}\to\mathcal{C}\times\mathcal{C}$ le foncteur
diagonal $\Delta(A)=(A,A)$. Si $\mathcal{C}$ admet des produits binaires,
alors $\Delta\dashv\times$~:
\[
    \Hom_{\mathcal{C}\times\mathcal{C}}(\Delta A,\,(B,C))
    = \Hom(A,B)\times\Hom(A,C)
    \cong \Hom(A,\,B\times C).
\]
Dualement, le coproduit est adjoint à gauche du diagonal~:
$\sqcup\dashv\Delta$.
\end{exemple}

\begin{exemple}[Extension et restriction des scalaires]\label{ex:scalaires}
\index{adjonction!extension-restriction des scalaires}
\index{extension des scalaires}
Soit $\varphi\colon R\to S$ un homomorphisme d'anneaux.
Le foncteur de restriction des scalaires $\varphi^*\colon\Mod_S\to\Mod_R$
admet un adjoint à gauche, l'extension des scalaires
$S\otimes_R(-)\colon\Mod_R\to\Mod_S$~:
\[
    \Hom_S(S\otimes_R M,\, N) \cong \Hom_R(M,\,\varphi^* N).
\]
\end{exemple}


%% ============================================================
\section{Adjoints et limites}\label{sec:adjoints-limites}
%% ============================================================

\begin{theoreme}[Les adjoints à droite préservent les limites]\label{thm:rapl}
\index{adjoint à droite!préserve les limites}
\index{RAPL}
Soit $F\dashv G$ avec $F\colon\mathcal{C}\to\mathcal{D}$ et
$G\colon\mathcal{D}\to\mathcal{C}$.
Alors $G$ préserve toutes les limites qui existent dans~$\mathcal{D}$.
\end{theoreme}

\begin{proof}
Soit $D\colon\mathcal{J}\to\mathcal{D}$ un diagramme avec limite
$(L,\pi_j)$. Montrons que $(GL,G\pi_j)$ est une limite de
$GD\colon\mathcal{J}\to\mathcal{C}$.

Soit $(X,\psi_j)$ un cône sur~$GD$ dans~$\mathcal{C}$,
i.e.\ $\psi_j\colon X\to GD_j$ avec $GD_\alpha\circ\psi_j=\psi_k$
pour tout $\alpha\colon j\to k$.

Par l'adjonction, chaque $\psi_j\colon X\to GD_j$ correspond à un unique
$\hat{\psi}_j = \varepsilon_{D_j}\circ F\psi_j\colon FX\to D_j$.
Vérifions que $(\hat{\psi}_j)_j$ est un cône sur~$D$.
Pour $\alpha\colon j\to k$~:
\begin{align*}
    D_\alpha\circ\hat{\psi}_j
    &= D_\alpha\circ\varepsilon_{D_j}\circ F\psi_j
    = \varepsilon_{D_k}\circ FGD_\alpha\circ F\psi_j \\
    &= \varepsilon_{D_k}\circ F(GD_\alpha\circ\psi_j)
    = \varepsilon_{D_k}\circ F\psi_k = \hat{\psi}_k,
\end{align*}
où l'on a utilisé la naturalité de~$\varepsilon$.

Par la propriété universelle de~$(L,\pi_j)$, il existe un unique
$v\colon FX\to L$ avec $\pi_j\circ v = \hat{\psi}_j$.
En transposant, on obtient $u = \tilde{v} = Gv\circ\eta_X\colon X\to GL$.
Alors~:
\[
    G\pi_j\circ u = G\pi_j\circ Gv\circ\eta_X
    = G(\pi_j\circ v)\circ\eta_X
    = G\hat{\psi}_j\circ\eta_X
    = \widetilde{\hat{\psi}}_j = \psi_j.
\]
L'unicité de~$u$ découle de celle de~$v$ et de la bijectivité de la transposition.
\end{proof}

\begin{theoreme}[Les adjoints à gauche préservent les colimites]\label{thm:lapc}
\index{adjoint à gauche!préserve les colimites}
\index{LAPC}
Soit $F\dashv G$ avec $F\colon\mathcal{C}\to\mathcal{D}$ et
$G\colon\mathcal{D}\to\mathcal{C}$.
Alors $F$ préserve toutes les colimites qui existent dans~$\mathcal{C}$.
\end{theoreme}

\begin{proof}
Soit $D\colon\mathcal{J}\to\mathcal{C}$ un diagramme avec colimite
$(C,\iota_j)$. Montrons que $(FC,F\iota_j)$ est une colimite de
$FD\colon\mathcal{J}\to\mathcal{D}$.

Soit $(Y,\psi_j)$ un cocône sous~$FD$, i.e.\ $\psi_j\colon FD_j\to Y$.
Par l'adjonction, chaque $\psi_j$ correspond à
$\tilde{\psi}_j\colon D_j\to GY$.
La famille $(\tilde{\psi}_j)_j$ est un cocône sous~$D$ dans~$\mathcal{C}$
(par naturalité de la bijection).

Par la propriété universelle de~$(C,\iota_j)$, il existe un unique
$w\colon C\to GY$ avec $w\circ\iota_j = \tilde{\psi}_j$.
En transposant, $\hat{w} = \varepsilon_Y\circ Fw\colon FC\to Y$ satisfait~:
\[
    \hat{w}\circ F\iota_j
    = \varepsilon_Y\circ Fw\circ F\iota_j
    = \varepsilon_Y\circ F(w\circ\iota_j)
    = \varepsilon_Y\circ F\tilde{\psi}_j
    = \hat{\tilde{\psi}}_j = \psi_j.
\]
L'unicité suit comme dans la preuve précédente.
\end{proof}

\begin{corollaire}\label{cor:adj-hom-lim}
\index{adjonction!et Hom}
Si $F\dashv G$, alors pour tout diagramme $D\colon\mathcal{J}\to\mathcal{D}$
dont la limite existe~:
\[
    G(\varprojlim D_j) \cong \varprojlim G(D_j),
    \qquad
    F(\varinjlim D_j) \cong \varinjlim F(D_j).
\]
\end{corollaire}


%% ============================================================
\section{Le théorème de Freyd (GAFT)}\label{sec:freyd}
%% ============================================================

Nous établissons maintenant un critère puissant garantissant l'existence
d'un foncteur adjoint.

\begin{definition}[Condition d'ensemble de solutions]\label{def:solution-set}
\index{condition d'ensemble de solutions}
\index{solution set condition}
Un foncteur $G\colon\mathcal{D}\to\mathcal{C}$ satisfait la
\textbf{condition d'ensemble de solutions} si, pour tout $A\in\mathcal{C}$,
il existe un ensemble (et non une classe propre)
$\{f_i\colon A\to GB_i\}_{i\in I_A}$ de morphismes tel que tout morphisme
$f\colon A\to GB$ se factorise comme $f = Gh\circ f_i$ pour un certain
$i\in I_A$ et un certain $h\colon B_i\to B$.
\[
\begin{tikzcd}
A \arrow[r, "f_i"] \arrow[dr, "f"'] & GB_i \arrow[d, "Gh"] \\
& GB
\end{tikzcd}
\]
\end{definition}

\begin{theoreme}[Théorème de l'adjoint de Freyd (GAFT)]\label{thm:gaft}
\index{théorème de Freyd}
\index{GAFT}
\index{théorème de l'adjoint}
Soit $G\colon\mathcal{D}\to\mathcal{C}$ un foncteur où $\mathcal{D}$
est complète (admet toutes les petites limites).
Alors $G$ admet un adjoint à gauche si et seulement si~:
\begin{enumerate}[label=(\roman*)]
\item $G$ préserve toutes les petites limites (est continu) ;
\item $G$ satisfait la condition d'ensemble de solutions.
\end{enumerate}
\end{theoreme}

\begin{proof}
\emph{Nécessité}~: si $F\dashv G$, alors $G$ préserve les limites par
le \autoref{thm:rapl}. Pour la condition d'ensemble de solutions,
l'ensemble $\{\eta_A\colon A\to GFA\}$ (réduit à un seul élément)
convient, car tout $f\colon A\to GB$ a pour transposée $\hat{f}\colon FA\to B$,
et $f = G\hat{f}\circ\eta_A$.

\emph{Suffisance}~: fixons $A\in\mathcal{C}$.
Nous construisons un objet initial dans la catégorie comma $(A\downarrow G)$\index{catégorie comma},
dont les objets sont les paires $(B,f\colon A\to GB)$ et les morphismes
$(B,f)\to(B',f')$ sont les $h\colon B\to B'$ avec $Gh\circ f = f'$.

\emph{Étape 1.} Par la condition d'ensemble de solutions, soit
$\{(B_i,f_i)\}_{i\in I}$ un ensemble de solutions.
Posons $B_0 = \prod_{i\in I}B_i$ (produit dans~$\mathcal{D}$,
qui existe car $\mathcal{D}$ est complète).
Comme $G$ préserve les limites, $GB_0\cong\prod_i GB_i$,
et la famille $(f_i)_i$ induit $f_0\colon A\to GB_0$.
L'objet $(B_0,f_0)$ est \og{}faiblement initial\fg{}~: tout objet
de $(A\downarrow G)$ reçoit au moins un morphisme depuis $(B_0,f_0)$.

\emph{Étape 2.} Pour rendre l'objet initial, nous égalisons.
Considérons l'ensemble de tous les endomorphismes
$h\colon B_0\to B_0$ dans $(A\downarrow G)$
(c'est un ensemble car $\Hom(B_0,B_0)$ est un ensemble).
Formons l'égaliseur $e\colon E\to B_0$ de tous ces endomorphismes
dans~$\mathcal{D}$.
(Plus précisément, l'égaliseur de la famille de paires
$(\id_{B_0}, h)$ pour tout endomorphisme~$h$.)
Comme $G$ préserve les égaliseurs, $Ge\colon GE\to GB_0$ est
l'égaliseur correspondant dans~$\mathcal{C}$.
Puisque $\id_{B_0}$ est un endomorphisme de $(B_0,f_0)$,
on a $Gh\circ f_0 = f_0$ pour tout~$h$, donc
$f_0$ se factorise par $Ge$ en un unique $f_E\colon A\to GE$.

\emph{Étape 3.} Montrons que $(E,f_E)$ est initial dans $(A\downarrow G)$.
Par l'étape~1 (via la factorisation $e$), tout objet $(B,f)$ reçoit un
morphisme depuis $(E,f_E)$~: en effet, le morphisme $(B_0,f_0)\to(B,f)$
se factorise par~$E$.
Pour l'unicité, soient $u,v\colon E\to B$ deux morphismes avec
$Gu\circ f_E = Gv\circ f_E = f$.
Formons l'égaliseur $w\colon W\to E$ de $u$ et~$v$.
Par faible initialité de~$(B_0,f_0)$, il existe $s\colon B_0\to W$.
La composée $e' = e\circ w\circ s\colon B_0\to B_0$
est un endomorphisme de $(B_0,f_0)$, donc $e\circ w\circ s$ est égalisé
par~$e$, ce qui implique $w\circ s\circ e_0 = \id$ (par la propriété
de l'égaliseur en étape~2), et finalement $w$ est un isomorphisme,
donc $u=v$.

Posons alors $FA = E$. La correspondance $A\mapsto FA$ s'étend
en un foncteur adjoint à gauche de~$G$, avec unité $\eta_A = f_E$.
\end{proof}

\begin{corollaire}\label{cor:saft}
\index{SAFT}
\index{théorème de l'adjoint spécial}
\textbf{(Théorème de l'adjoint spécial, SAFT.)}
Soit $G\colon\mathcal{D}\to\mathcal{C}$ un foncteur continu
où $\mathcal{D}$ est complète et localement petite.
Si $\mathcal{D}$ admet une famille cogénératrice\index{famille cogénératrice},
alors $G$ admet un adjoint à gauche.
\end{corollaire}

\begin{proof}[Esquisse]
La famille cogénératrice fournit automatiquement la condition d'ensemble
de solutions, et on applique le \autoref{thm:gaft}.
\end{proof}

\begin{exemple}\label{ex:applications-gaft}
\begin{enumerate}[label=(\roman*)]
\item \index{complétion!de groupe abélien}
      Le foncteur d'oubli $U\colon\Ab\to\Set$ préserve les limites
      et $\Ab$ est complète. L'ensemble $\{f\colon X\to UB : |B|\leq|X|+\aleph_0\}$
      fournit un ensemble de solutions, d'où $F\dashv U$ (groupes abéliens libres).
\item \index{théorème de Stone-Čech}
      Le foncteur d'oubli $U\colon\mathsf{CompHaus}\to\Top$
      satisfait les conditions du SAFT, ce qui donne la
      compactification de Stone--Čech $\beta\dashv U$.
\end{enumerate}
\end{exemple}


%% ============================================================
\section{Adjonctions et catégories comma}\label{sec:adj-comma}
%% ============================================================

\begin{proposition}\label{prop:adj-objet-initial-comma}
\index{adjonction!et catégorie comma}
\index{catégorie comma}
Un foncteur $G\colon\mathcal{D}\to\mathcal{C}$ admet un adjoint à gauche
si et seulement si, pour tout $A\in\mathcal{C}$, la catégorie comma
$(A\downarrow G)$ admet un objet initial.
\end{proposition}

\begin{proof}
Si $F\dashv G$ avec unité~$\eta$, alors $(FA,\eta_A)$ est initial
dans $(A\downarrow G)$~: pour tout $(B,f\colon A\to GB)$, le morphisme
$\hat{f}=\varepsilon_B\circ Ff\colon FA\to B$ est l'unique morphisme
avec $G\hat{f}\circ\eta_A = f$.

Réciproquement, si $(L_A,\eta_A\colon A\to GL_A)$ est initial dans
$(A\downarrow G)$ pour tout~$A$, on pose $FA=L_A$ et on vérifie
la fonctorialité et la bijection naturelle.
\end{proof}


%% ============================================================
\section{Propriétés supplémentaires}\label{sec:adj-proprietes}
%% ============================================================

\begin{proposition}\label{prop:adj-pleinement-fidele}
\index{adjonction!et foncteur pleinement fidèle}
\index{foncteur pleinement fidèle!et adjonction}
Soit $F\dashv G$.
\begin{enumerate}[label=(\roman*)]
\item $G$ est fidèle si et seulement si chaque composante de la
      counité~$\varepsilon_B$ est un épimorphisme.
\item $G$ est pleinement fidèle si et seulement si $\varepsilon$
      est un isomorphisme naturel.
\end{enumerate}
\end{proposition}

\begin{proof}
(i) $G$ fidèle $\Leftrightarrow$ $\Phi$ injective sur les fibres
$\Leftrightarrow$ la composée $h\mapsto\varepsilon_B\circ Fk\mapsto h$
est injective
$\Leftrightarrow$ $\varepsilon_B$ est un épimorphisme (car $\hat{k}=\varepsilon_B\circ Fk$).

(ii) $G$ pleinement fidèle $\Leftrightarrow$ $\Phi$ bijectif pour $A=GB$
$\Leftrightarrow$ $\Hom(FGB,B)\cong\Hom(GB,GB)$ par $h\mapsto Gh\circ\eta_{GB}$,
et $\varepsilon_B$ est l'antécédent de $\id_{GB}$, d'où le résultat.
\end{proof}

\begin{proposition}[Unicité des adjoints]\label{prop:unicite-adjoints}
\index{adjonction!unicité}
Si $F\dashv G$ et $F'\dashv G$, alors $F\cong F'$ (naturellement).
De même, si $F\dashv G$ et $F\dashv G'$, alors $G\cong G'$.
\end{proposition}

\begin{proof}
Par le lemme de Yoneda, $F$ est déterminé par le foncteur
$\Hom(F(-),-)=\Hom(-,G(-))$, qui ne dépend que de~$G$.
\end{proof}

\begin{proposition}[Composition d'adjonctions]\label{prop:composition-adj}
\index{adjonction!composition}
Si $F\dashv G\colon\mathcal{C}\rightleftarrows\mathcal{D}$ et
$F'\dashv G'\colon\mathcal{D}\rightleftarrows\mathcal{E}$, alors
$F'F\dashv GG'$.
\end{proposition}

\begin{proof}
$\Hom_\mathcal{E}(F'FA,B)
\cong\Hom_\mathcal{D}(FA,G'B)
\cong\Hom_\mathcal{C}(A,GG'B)$,
naturellement en $A$ et~$B$.
\end{proof}


%% ============================================================
\section{Monades issues d'adjonctions}\label{sec:monades-apercu}
%% ============================================================

\begin{remarque}\label{rem:monade-apercu}
\index{monade}
Toute adjonction $F\dashv G$ engendre une \emph{monade}
$T = GF\colon\mathcal{C}\to\mathcal{C}$ munie de~:
\begin{itemize}
\item $\eta\colon\Id\Rightarrow T$ (l'unité de l'adjonction),
\item $\mu = G\varepsilon F\colon T^2 = GFGF\Rightarrow GF = T$
      (la multiplication).
\end{itemize}
Les identités triangulaires entraînent les axiomes de monade~:
$\mu\circ T\eta = \mu\circ\eta T = \id_T$ et
$\mu\circ T\mu = \mu\circ\mu T$ (associativité).
Nous développerons cette théorie au chapitre suivant.
\end{remarque}


%% ============================================================
\section{Exercices}\label{sec:exo-adjonctions}
%% ============================================================

\begin{exercice}\label{exo:adj-libre-ab}
\index{adjonction!groupe abélien libre}
Soit $F\colon\Set\to\Ab$ le foncteur \og{}groupe abélien libre\fg{}
($FX = \bigoplus_{x\in X}\Z$) et $U\colon\Ab\to\Set$ l'oubli.
Montrer que $F\dashv U$ en exhibant l'unité, la counité,
et en vérifiant les identités triangulaires.
\end{exercice}

\begin{exercice}\label{exo:adj-tensor-hom}
\index{adjonction!tenseur-hom}
Soit $R$ un anneau commutatif et $M$ un $R$-module.
Montrer que $M\otimes_R(-)\dashv\Hom_R(M,-)$ en établissant
la bijection naturelle
$\Hom_R(M\otimes_R N,\,P)\cong\Hom_R(N,\,\Hom_R(M,P))$.
\end{exercice}

\begin{exercice}\label{exo:adj-preserve-epi}
\index{adjoint à gauche!préserve les épimorphismes}
Montrer que tout adjoint à gauche préserve les épimorphismes
et que tout adjoint à droite préserve les monomorphismes.
\end{exercice}

\begin{exercice}\label{exo:adj-equivalence}
\index{équivalence de catégories!et adjonction}
Montrer qu'une adjonction $F\dashv G$ est une équivalence de catégories
si et seulement si $\eta$ et $\varepsilon$ sont des isomorphismes naturels.
\end{exercice}

\begin{exercice}\label{exo:adj-reflexive}
\index{sous-catégorie réflexive}
Soit $\mathcal{D}\hookrightarrow\mathcal{C}$ une sous-catégorie pleine.
On dit que $\mathcal{D}$ est \emph{réflexive} si le foncteur d'inclusion
$\iota\colon\mathcal{D}\hookrightarrow\mathcal{C}$ admet un adjoint à gauche
$L\colon\mathcal{C}\to\mathcal{D}$ (le \emph{réflecteur}).
\begin{enumerate}[label=(\alph*)]
\item Montrer que $\varepsilon$ est un isomorphisme naturel.
\item Montrer que $L$ est essentiellement surjectif.
\item Montrer que $\Ab$ est une sous-catégorie réflexive de $\Grp$
      (le réflecteur est l'abélianisation $G\mapsto G/[G,G]$).
\end{enumerate}
\end{exercice}

\begin{exercice}\label{exo:gaft-application}
\index{théorème de Freyd!application}
En utilisant le GAFT, montrer que le foncteur d'oubli
$U\colon\mathsf{CompHaus}\to\Top$ admet un adjoint à gauche
(la compactification de Stone--Čech).
\emph{Indication}~: vérifier que $\mathsf{CompHaus}$ est complète,
que $U$ préserve les limites, et exhiber un ensemble de solutions.
\end{exercice}

\begin{exercice}\label{exo:adj-kan}
\index{extension de Kan!à gauche}
Soient $K\colon\mathcal{A}\to\mathcal{B}$ et $F\colon\mathcal{A}\to\mathcal{C}$
des foncteurs avec $\mathcal{A}$ petite et $\mathcal{C}$ cocomplète.
Montrer que l'extension de Kan à gauche $\mathrm{Lan}_K F$
existe et que $\mathrm{Lan}_K\dashv K^*$ où
$K^*\colon\Fun(\mathcal{B},\mathcal{C})\to\Fun(\mathcal{A},\mathcal{C})$
est la précomposition par~$K$.
\end{exercice}

\begin{exercice}\label{exo:adj-monade}
\index{monade!axiomes}
Soit $F\dashv G$ avec $T=GF$, $\eta$ l'unité et $\mu=G\varepsilon F$.
Vérifier en détail les axiomes de monade~:
\begin{enumerate}[label=(\alph*)]
\item $\mu\circ T\eta = \id_T$ et $\mu\circ\eta T = \id_T$
      (axiomes d'unité).
\item $\mu\circ T\mu = \mu\circ\mu T$ (associativité).
\end{enumerate}
\emph{Indication}~: utiliser les identités triangulaires et la
naturalité de $\eta$ et $\varepsilon$.
\end{exercice}

\begin{exercice}\label{exo:adj-galois}
\index{correspondance de Galois}
\index{adjonction!de Galois}
Soient $(P,\leq)$ et $(Q,\leq)$ deux ensembles partiellement ordonnés,
vus comme catégories. Montrer qu'une adjonction $F\dashv G$
entre ces catégories est exactement une \emph{connexion de Galois}~:
$F(p)\leq q \Leftrightarrow p\leq G(q)$.
Donner un exemple en théorie des treillis.
\end{exercice}

\index{adjonction|)}
